\section*{第 \ref{ch:b1:operational-amplifier} 章——运算放大器}

\begin{solution}{exo:b1:operational-amplifier:1}
$G = 1 + 9 = 10$；$V_s = \qty{5.0}{V}$；饱和发生在 $V_e = 15/10 =
\qty{1.5}{V}$ 处。
\end{solution}

\begin{solution}{exo:b1:operational-amplifier:2}
$G = -10$；输入阻抗 $R_1 = \qty{10}{k\ohm}$；$i = 1.0/10^4 =
\qty{0.10}{mA}$；$V_s = \qty{-10}{V}$。
\end{solution}

\begin{solution}{exo:b1:operational-amplifier:3}
不接跟随器时，分压比 $1/(1 + 10) = 0.091$，即 \qty{91}{mV}。接入跟随器后，
输出为 \qty{1.0}{V}——跟随器不从传感器取电流，
并提供负载所需的 \qty{1}{mA}。
\end{solution}

\begin{solution}{exo:b1:operational-amplifier:4}
$V_s = -R_f(V_1/R_1 + V_2/R_2) = -(V_1 + V_2/2) = -(1.0 + 1.0) = \qty{-2.0}{V}$。
\end{solution}

\begin{solution}{exo:b1:operational-amplifier:5}
$V_+ = V_2R_2/(R_1 + R_2)$；在 $V_-$ 处：$(V_1 - V_-)/R_1 = (V_- - V_s)/R_2$，
故 $V_s = V_-(1 + R_2/R_1) - V_1R_2/R_1$；由 $V_- = V_+$ 可得
$V_s = V_2R_2/R_1 - V_1R_2/R_1$。代入数值：$10 \times 0.05 = \qty{0.50}{V}$；
公共的 \qty{1}{V} 不产生任何贡献——这正是共模抑制。
\end{solution}

\begin{solution}{exo:b1:operational-amplifier:6}
$V_s = -(1/RC)\int V_e$：斜率 $-1000 \times 1.0 = \qty{-1000}{V/s}$；
经过 $15/1000 = \qty{15}{ms}$ 后达到饱和。对方波输入，输出为三角波，
振幅为 $ET/4RC = 1.0 \times 10^{-3}/(4 \times 10^{-3}) = \qty{0.25}{V}$
（再加上初始电荷留下的任何常数）。
\end{solution}

\begin{solution}{exo:b1:operational-amplifier:7}
$\underline H = -(R'/R)/(1 + jR'C\omega)$：直流增益 $-100$；截止频率
$1/(2\pi R'C) = \qty{1.6}{Hz}$；在 \qty{1}{kHz} 处，$R'C\omega = 628$，$G =
100/628 = 0.16$——即积分器的值 $1/RC\omega$。
当 $f \gg \qty{1.6}{Hz}$ 时电路起积分作用。
\end{solution}

\begin{solution}{exo:b1:operational-amplifier:8}
$R_2/R_1 = 10$：$R_2 = \qty{10}{k\ohm}$；$C = 1/(2\pi R_2f_c) = 1/(2\pi
\times 10^4 \times 500) = \qty{32}{nF}$。
\end{solution}

\begin{solution}{exo:b1:operational-amplifier:9}
考虑有限的 $\mu$：$V_s = \mu(V_e - \beta V_s)$，$\beta = 1/100$，故
$G = \mu/(1 + \mu\beta) = 10^5/(1 + 10^3) = 99.9$。带宽为 $f_T/G =
\qty{10}{kHz}$；跟随器的带宽为 \qty{1}{MHz}。
\end{solution}

\begin{solution}{exo:b1:operational-amplifier:10}
噪声每次越过 \qty{2.0}{V} 都会使输出翻转，在过零点附近出现一串快速开关
（即“抖动”）。补救办法是采用比噪声更宽的滞回，
即施密特触发器。
\end{solution}

\begin{solution}{exo:b1:operational-amplifier:11}
$V_+ = V_sR_1/(R_1 + R_2) = V_s/10$：阈值为 $\pm\qty{1.5}{V}$。
输出保持 $+\qty{15}{V}$，直到 $V_e$ 上升至 \qty{1.5}{V} 时翻转为
$-\qty{15}{V}$，并保持该状态，直到 $V_e$ 降至 \qty{-1.5}{V} 以下时翻转
回来：得到与三角波频率相同的方波，其跳变沿位于
阈值穿越处。
\end{solution}

\begin{solution}{exo:b1:operational-amplifier:12}
当 $V_s = +V_{\mathrm{sat}}$ 时，$V_C$ 从 $-\beta V_{\mathrm{sat}}$
向 $+V_{\mathrm{sat}}$ 充电：$V_C = V_{\mathrm{sat}} - (1 + \beta)V_{\mathrm{sat}}
\eu^{-t/RC}$，达到 $+\beta V_{\mathrm{sat}}$ 时，$\eu^{-t/RC} =
(1 - \beta)/(1 + \beta)$：半周期为 $RC\ln\frac{1 + \beta}{1 - \beta}$，
由此得到 $T$。取 $\beta = \tfrac12$：$T = 2RC\ln 3 = 2.2RC$；$RC = 1/(2.2 \times
10^3) = \qty{0.455}{ms}$，$R = \qty{4.6}{k\ohm}$。$V_s$ 为方波
$\pm\qty{15}{V}$；$V_C$ 为在 $\pm\qty{7.5}{V}$ 之间变化的指数弧线。
\end{solution}

\begin{solution}{pb:b1:operational-amplifier:1}
\textbf{1.} 两个分压器：$u = E[\tfrac12 - (1 + \epsilon)/(2 + \epsilon)]
\approx -E\epsilon/4$（见\cref{prop:b1:dc-circuits:wheatstone}）；
接线方式决定符号。

\textbf{2.} $u = 5.00 \times 2.0\times10^{-3}/4 = \qty{2.5}{mV}$。

\textbf{3.} $V_A = E/2 = \qty{2.500}{V}$ 始终成立；满量程时 $V_B = E(1 + \epsilon)/(2 +
\epsilon) \approx 2.500 + 0.0025 = \qty{2.5025}{V}$。
两者共有的 \qty{2.5}{V} 即为共模电压。

\textbf{4.} 每个中点：$1 \parallel 1 = \qty{500}{\ohm}$；
两中点之间为 \qty{1.0}{k\ohm}。

\textbf{5.} $10/(10 + 1) = 0.91$：损失 $9\%$。

\textbf{6.} 输入阻抗无穷大（电桥不提供电流，
故中点保持开路时的电势），输出阻抗
为零（下一级可取用所需的任何电流）。

\textbf{7.} 恰为 $V_A$ 和 $V_B$；电流为零。

\textbf{8.} 与\cref{prop:b1:operational-amplifier:sumdiff}中相同：$V_+ =
V_BR_2/(R_1 + R_2)$，在 $V_-$ 处用节点定律，并取 $V_- = V_+$。

\textbf{9.} $R_2 = \qty{1.0}{M\ohm}$。

\textbf{10.} $1 + R_2/R_1 = 10$：例如 \qty{1}{k\ohm} 与 \qty{9}{k\ohm}。
总增益为 $1000$；满量程为 $2.5 \times 10^{-3} \times 1000 = \qty{2.5}{V}$。

\textbf{11.} 匹配时，公共的 \qty{2.5}{V} 被精确抵消（仅剩 $V_B - V_A$）。
$1\%$ 的失配会使其中约 $1\%$ 通过：
折合到输入端的误差量级为 \qty{25}{mV}，是满量程信号的十倍——
四个电阻的匹配精度必须优于 $0.01\%$
（或者先去除共模电压）。

\textbf{12.} $1\ \unit{mV} \times 1000 = \qty{1}{V}$，占满量程的 40\%：
若不做失调调零，不用低失调（斩波式）
放大器，或在梁上无负载时做零点校准，则不可用。

\textbf{13.} 第一级：$f_T/100 = \qty{10}{kHz}$；第二级：\qty{100}{kHz}；
均远高于 \qty{5}{Hz}。

\textbf{14.} 所需压摆率为 $2.5/10^{-2} = \qty{250}{V/s} = \qty{0.25}{mV/\micro s}$，
比压摆率低两千倍：否。

\textbf{15.} $C = 1/(2\pi R_2f_c) = 1/(2\pi \times 1.6\times10^4 \times 10)
= \qty{1.0}{\micro F}$.

\textbf{16.} $x = 5$: $G = 1/\sqrt{26}$, \qty{-14}{dB}.

\textbf{17.} $\varphi = -\arctan 0.5 = -27^\circ$；延迟为 $\varphi/\omega =
\qty{15}{ms}$——与卡车驶过相比可忽略不计。

\textbf{18.} 对后级无负载效应（输出阻抗为零）
且可提供增益；同时，有用信号在低于截止
频率时不会被衰减。

\textbf{19.} \qty{-28}{dB}（增益相乘，分贝相加），精确成立，因为
第一级的输出阻抗为零：第二级不会
对其造成负载效应。

\textbf{20.} $\epsilon = 1.6\times10^{-3} \to u = \qty{2.0}{mV} \to
\qty{2.0}{V}$；输出为 $+V_{\mathrm{sat}}$，当信号超过
\qty{2.0}{V}时，低于时为 $-V_{\mathrm{sat}}$（或因接线相反而互换）。

\textbf{21.} 会颤动：噪声穿过参考
电压时，灯会闪烁。

\textbf{22.} 将 $V_{\mathrm{ref}}$ 经 $R_2$ 接入，输出经
$R_1$\dots{} 按标准形式接回，同相输入端电位为
$V_{\mathrm{ref}}(1 - \beta) + \beta V_s$；相应地取平移后的参考电压，
则翻转电平为 $V_{\mathrm{ref}} \pm \beta V_{\mathrm{sat}}$。
$\beta V_{\mathrm{sat}} = \qty{0.10}{V}$：$\beta = 1/150$，例如 $R_1 =
\qty{1.0}{k\ohm}$，$R_2 = \qty{149}{k\ohm}$。

\textbf{23.} \qty{2.1}{V}（触发）与 \qty{1.9}{V}（释放）：应变分别为
\num{1.68e-3} 与 \num{1.52e-3}。

\textbf{24.} $R = (15 - 2.0)/0.013 = \qty{1.0}{k\ohm}$。
在 $-V_{\mathrm{sat}}$ 时，发光二极管承受 \qty{15}{V} 的反向偏压，超出其
额定值：反并联（或串联）一个普通二极管即可保护它。

\textbf{25.} 电桥（满量程 \qty{2.5}{mV}）$\to$ 两个跟随器（无
负载效应）$\to$ 差分放大器 $\times100$ $\to$ 同相
放大器 $\times10$ $\to$ 有源低通滤波器 \qty{10}{Hz} $\to$ 施密特触发器，
阈值为 $\qty{2.0}{V} \pm \qty{0.1}{V}$ $\to$ 发光二极管。真正的限制在于：输入
失调电压，以及差分放大器各电阻的匹配，
两者都与毫伏量级的信号相当。
\end{solution}
